% REPORT-001 中文审阅稿。译自按核验意见修订后的 main.tex（2026-10-28），仅供作者审阅、不投稿；权威稿仍为 main.tex。
% 构建：在 b-Office/technical-report 下执行两遍
%   xelatex -interaction=nonstopmode -halt-on-error -output-directory=tmp/zh -jobname=echo-technical-report-zh main-zh.tex
% （需 ctex 与 Windows 中文字体；tmp/ 已忽略入库）。蓝色【改 n】标出按同目录 REPORT-001-核验-2026-10-28.md 第 n 条修改之处。
\documentclass[11pt,a4paper,fontset=windows]{ctexart}
\usepackage[margin=23mm]{geometry}
\usepackage{amsmath,booktabs,array,graphicx,xcolor}
\usepackage{tikz}
\usetikzlibrary{arrows.meta,positioning}
\usepackage[hidelinks]{hyperref}
\pagestyle{plain}
\setlength{\parskip}{3pt}
\setlength{\emergencystretch}{1.5em}
\newcommand{\code}[1]{\texttt{#1}}
\newcommand{\audit}[1]{\textcolor{blue}{\textsuperscript{【改#1】}}}
\title{\vspace{-1.1cm}\textbf{Echo：通过逐词引导对话模型\\来组织第二语言回复}\vspace{-0.2cm}}
\author{Yuxuan Chen\\
\small 宁波工程学院（Ningbo University of Technology）\quad $\cdot$ \quad \texttt{cyx3705@outlook.com}}
\date{\vspace{-0.8em}\small 技术报告 \quad $\cdot$ \quad 2026 年 10 月}
\begin{document}
\maketitle
\thispagestyle{empty}

\noindent\fbox{\parbox{\dimexpr\linewidth-2\fboxsep-2\fboxrule}{\small
\textbf{审阅说明。}本稿逐段译自按核验意见修订后的 \code{main.tex}（2026-10-28），数字与英文稿一致。蓝色【改 n】标出按 \code{REPORT-001-核验-2026-10-28.md} 第 n 条修改的位置，文末附各条处理情况。第 5 节全部数字已于 2026-10-28 在提交 \code{143bec4} 上重新测量，原始输出见 \code{data/2026-10-28/}；测试状态、终端环境等运维记录已删除。英文 PDF 已用修订稿重新生成\audit{1}。}}

\begin{abstract}
第二语言学习者能认出的词往往远多于能主动说出的词，而 AI 写作辅助通常走两个极端：要么让学习者独自面对空白输入框，要么直接替他们写好回复。Echo 是一个面向以中文为母语的英语学习者的本地原型，探索两者之间的做法：由学习者充当解码器。扮演学习者对话伙伴的同一个对话模型，被用来续写学习者自己尚未完成的这一轮；界面把它的下一词分布呈现为一小列带释义的候选词。学习者要么接受其中一个，要么自由输入，于是回复沿着模型的续写树、由学习者一步步选出，任何文字都必须经过学习者的操作才会进入消息。这个循环也覆盖学习者暂时说不出的部分：就地输入的中文词，会被替换为按语境选出的英文词典候选。它同样覆盖学习者要读的部分：伙伴的回复带有逐词中文释义，释义从词典义项中选出；学习者把某个词元标为已会后释义即隐藏，系统退回默认释义时会作可见标记。这个循环只有在候选先于下一次击键到达时才可用。因此 Echo 在学习者未结束的回合上做一次前向读取下一词排名，从不把半截词送入模型，并复用上一次击键的 KV 状态。在一块笔记本 GPU 上，每条输入消息除首次请求外，每次击键都在 100\,ms 内得到回应（中位 5\,ms；需要前向时约 20\,ms）\audit{2}\audit{4}；对 94\% 的击键，学习者实际打出的词就在五个候选之中\audit{8}。上下文变化后的首次请求约 220\,ms。一种直截了当的提示词工程替代方案，是通过 API 向同一模型询问「学习者最可能接着输入的 5 个词」。它每次击键的中位耗时为 358\,ms，从未在打字者平均击键间隔 239\,ms 内返回，5 个词中平均只有 1.7 个合格候选，只有 48\% 的击键包含实际打出的词。同一运行时把一条伙伴回复的全部词典受限释义显示上屏需 2.4\,s，经 API 则需 5.4\,s。本文尚未开展用户研究；我们给出设计所引出的假设，并概述检验方案。
\end{abstract}

\section{引言}
即使学习者认得一个完整英文句子所需的大部分词，自己说出这个句子仍可能很难。两种常见的 AI 支持位于两端。空白输入框加可选的语法检查，并不触及「产出」这一难题。起草或改写整条回复的系统消除了这个难题，也随之拿走了学习者自己的大部分产出。Echo 探索两者之间的空间。

在 Echo 中，学习者用英语与本地语言模型对话。学习者组句时，\emph{同一个}模型以相同的对话历史和指令，被要求续写学习者未完成的这一轮。它排序后的下一词预测以少量候选显示，每个候选附简短中文释义。学习者可以接受其中一个、继续输入或编辑。实际上，学习者取代了对话语言模型的采样器：每次接受都确定模型续写树中的一个分支，由一个有话要说的人来选，而不是随机抽取。当某个词完全想不出来时，学习者可以用中文输入，让 Echo 按语境选出对应的英文词典条目。阅读伙伴回复时，学习者会看到尚未标为已会的词的中文释义。

本报告有五项贡献：
\begin{itemize}\setlength{\itemsep}{1pt}
\item \textbf{一种交互范式}：在第二语言对话中，学习者逐词引导对话模型对学习者自己回合的预测，同时在同一对话中为伙伴回合提供理解支持（第~\ref{sec:paradigm}~节）。
\item \textbf{该循环的具体设计}及其理由：带释义的候选、半词处理、就地中译英查询、移除释义的词元级「已会」状态，以及不确定释义的可见标记（第~\ref{sec:design}~节）。
\item \textbf{单 GPU 实现}：通过共享同一模型已计算的上下文、并把模型输出限制为词典选项，使预测和选义保持跟手（第~\ref{sec:implementation}~节）。
\item \textbf{对照 HCI 阈值的延迟评测}：逐击键的下一输入预测，与通过 API 查询同一模型相比，分别采用相同的一次前向表述和提示词工程。结果表明，决定可用与否的是预测的表述方式；不把半截词送入模型、复用 KV 状态则进一步提供余量\audit{4}。释义与查询也做了对比，其中包括 API 更快的情形（第~\ref{sec:results}~节）。
\item \textbf{研究议程}：用户研究需要检验的假设（第~\ref{sec:limitations}~节）。
\end{itemize}
模型引导的文本输入、预测键盘和人机协同写作都有成熟的先例（第~\ref{sec:related}~节）。我们不主张首创。我们的贡献在于把三者结合进同一个对话循环：学习者引导解码自己的回合、第二语言场景、理解支持。

\section{交互范式：学习者充当解码器}
\label{sec:paradigm}
设 $h$ 为目前为止的对话，$s$ 为伙伴的系统指令，$u_{<t}$ 为学习者尚未完成的消息。Echo 用模型自身的对话模板渲染 $h$，并让学习者的回合保持\emph{开放}、止于 $u_{<t}$，于是模型的下一 token 分布 $p_\theta(\cdot\mid s,h,u_{<t})$ 就是对学习者接下来要说什么的预测。若直接从这个分布生成，就等于让模型替学习者说话。Echo 改为展示它经过过滤的前 $k$ 项，由学习者选择：
\begin{equation}
 u_t \in \mathrm{Filter}_k\!\left(\operatorname{top20} p_\theta(\cdot\mid s,h,u_{<t})\right)\;\cup\;\{\text{学习者输入的任何内容}\}.
\end{equation}
这里 $p_\theta$ 定义在 token 上，而 $u_t$ 是一个完整的词：$\mathrm{Filter}_k$ 只保留带前导空格、能构成词典词的 token；学习者正在打一个词时，上下文止于该词之前（第~\ref{sec:design}~节）\audit{18}。候选来自伙伴自己的上下文，因而反映话题和伙伴刚提出的问题。每次接受都会追加一个词并立即请求下一次排名，于是一条回复就是模型续写树中的一条路径，由学习者一层一层走出来。Echo 并不显示这棵树本身，只显示当前一层的备选，这让选择保持很小，也让学习者向前阅读，而不是浏览各种补全。

\begin{figure}[ht]
\centering
\begin{tikzpicture}[font=\small,>=Stealth,
 ctx/.style={draw,rounded corners=2pt,fill=black!5,align=left,inner sep=4pt},
 cand/.style={draw,rounded corners=2pt,inner sep=3pt,minimum width=15mm,align=center},
 pick/.style={cand,very thick,fill=black!10}]
\node[ctx] (c1) at (0,0) {\code{I want to learn\textvisiblespace}};
\node[pick] (k11) at (4.2,0.9) {English};
\node[cand] (k12) at (4.2,0.3) {how};
\node[cand] (k13) at (4.2,-0.3) {to};
\node[cand] (k14) at (4.2,-0.9) {more};
\node[cand] (k15) at (4.2,-1.5) {about};
\foreach \k in {k11,k12,k13,k14,k15} \draw[->] (c1.east) -- (\k.west);
\node[ctx] (c2) at (8.6,0.9) {\code{I want to learn English\textvisiblespace}\\[-1pt]{\scriptsize 立即请求下一次排名}};
\draw[->,very thick] (k11.east) -- (c2.west);
\node[ctx] (c3) at (0,-3.1) {\code{I want to mak}};
\node[cand] (m1) at (4.2,-2.45) {make};
\node[cand,dashed] (m2) at (4.2,-3.1) {maker, makeup};
\node[cand,dashed] (m3) at (4.2,-3.75) {making, makeshift};
\foreach \k in {m1,m2,m3} \draw[->] (c3.east) -- (\k.west);
\node[align=left,font=\scriptsize,text width=52mm] at (8.6,-3.1) {半词：排名在 \emph{mak} 之前取得，再按前缀过滤。实线框：模型候选；虚线框：按词典词频补齐。};
\end{tikzpicture}
\caption{两份记录下来的候选列表（\code{qwen3:8b}，5 个候选；2026 年 10 月 23 日记录，10 月 28 日的运行结果完全相同）\audit{17}。上：学习者用 Shift+空格接受 \emph{English}（示意性的选择），随即请求下一层。每个显示的候选还带有简短的中文词典释义，此处略去。}
\label{fig:walk}
\end{figure}

有三个特性把它与最常见的 AI 写作支持区分开来。第一，\emph{没有学习者的操作，任何内容都不会进入消息}：不会自动补全整句或整条回复，学习者也可以完全忽略候选。第二，预测来自\emph{对话伙伴}的上下文，而不是通用键盘模型或单独的辅导模型。第三，每一步都是学习者可以审视的词汇决定，因为候选同时显示词和释义。预期效果是把产出转化为一连串的识别决定，同时让学习者始终掌控意义。这究竟有助于还是有碍于学习，是本报告没有回答的实证问题（第~\ref{sec:limitations}~节）。

\section{设计}
\label{sec:design}
设计遵循五条工作原则：
\textbf{（D1）}学习者自主：模型提出，学习者决定；
\textbf{（D2）}有界的模型输出：展示给学习者的中文来自词典，绝不来自自由生成；
\textbf{（D3）}如实呈现不确定性：退回的默认释义要可见地标出；
\textbf{（D4）}注意力经济：学习者已会的词不再获得支持；
\textbf{（D5）}跟手：候选必须在一块消费级 GPU 上跟上打字速度。

\subsection{输入框}
一次前向给出下一 token 的前 20 名排名。Echo 保留带前导空格、且在词典中有条目的英文词（单字母只保留 \emph{a} 和 \emph{I}）以及四个标点，按不区分大小写去重，显示至多所配置的数量（默认 5 个，可设 1--20）。每个候选显示序号、词和简短中文释义。Shift+W 和 Shift+S 移动选择，Shift+空格接受（Tab 和点击也可接受）。普通空格输入空格，Enter 发送，因此普通输入不会被拦截（D1），只有一个例外：候选列表有效时，Shift+W 和 Shift+S 用于选择，而不是输入大写 \emph{W} 或 \emph{S}。由于伙伴一说完就会出现首批候选，这会影响 \emph{What}、\emph{So} 等回合开头的词；换一组键位即可消除冲突\audit{3}。接受会追加该词和一个空格，并请求下一次排名。在回合开头只保留首字母大写的候选，因为回合开头的小写 token 大多是词的碎片。半截词从不送入模型，因为在 token 中间截断的提示会产生 \emph{makng} 这样的拼写错误。Echo 在该词之前做排名，按已输入前缀过滤，模型给出的不够时再按词典词频补齐（图~\ref{fig:walk}）。候选行在更新时保持固定高度，输入框不会跳动。词典成员资格无法证明某个 token 恰好结束一个词：多 token 的词可能漏掉，碎片也可能被显示。

\subsection{说出暂时还不会用英语说的话}
\label{sec:reverse}
按住 Ctrl 期间，嵌在输入框中的原生子控件承载系统中文输入法。松开 Ctrl 后，已确认的中文写回光标处。随后一次普通空格把前面连续的一段中文（可能跨越多次 Ctrl 输入）确认为一个词，长度 1--32 个字符。Echo 在英汉词典中反查（至多 20 个词或短语），让模型根据对话和该词周围的文字从中选择。选中的条目就地替换中文。若没有条目或模型弃选，中文保持不变（D2）：Echo 从不退回自由翻译。例如 苹果$\to$\emph{apple}、冰淇淋$\to$\emph{ice cream}；实测延迟见第~\ref{sec:results}~节。请求进行期间，预测和发送暂停，但仍可编辑。过期的结果绝不会覆盖之后的编辑。

\subsection{阅读伙伴的回合}
伙伴回复以英文流式输出，每个词下方有一个小号中文释义。只有一个词典义项的词直接给出释义。对于多义词，模型在该词的词典义项中选择或弃选，只依据到该词为止的对话，绝不使用伙伴尚未生成的文字。选择结果随回复的继续而流式到达。若选择弃选、超出有界队列、超时或失败，Echo 显示确定性的词典默认释义并加下划线，悬停可看原因（D3）。下划线标记的是已识别的退回，它并不能发现所有错误含义，因此自信但错误的释义可能不带标记。点击一个词即把它的词元标为已会（\emph{went} 标记 \emph{go}）。此后该词元的释义在各处隐藏但保留占位，版面不会重排，Echo 也完全停止为该词元计算释义（D4）。学习状态保存在本地，并与一个浏览器扩展共享，该扩展对网页文字应用同样的释义。在网页中，可见文字按约 600 字符分块提供，因此网页释义可能用到同一块中靠后的文字。同一网页或同一视频内，先前的分块按阅读顺序保留为上下文（至多 2,048 token），每次请求只预填新的分块。学习者转到别处时，Echo 把这段上下文的文字和 KV 状态存入本地数据目录（至多 16\,GB、128 个文件），返回时再恢复。

\section{让循环跟手}
\label{sec:implementation}
本节描述的是公开仓库（第~\ref{sec:reproduction}~节）提交 \code{143bec4} 时的 Echo，即第~\ref{sec:results}~节所测的代码（包版本 0.3.0）。另一条分支后来把这里描述的运行时换成了外部推理服务，但没有改动版本号，因此确定系统的是提交号，而不是版本号\audit{13}。在该提交中，Echo 是一个带浏览器界面的本地 FastAPI 服务。Python 负责 HTTP/SSE 传输、提示、词典访问、学习状态和编排。基于 llama.cpp~\cite{llamacpp} 的 C++ 运行时负责分词、KV 状态、批处理、采样、下一 token 排名、受限选择和分支清理。它在一个固定的 4,096 token KV 池（三个序列槽）中运行一个量化 8B 模型（\code{qwen3:8b}，Q4\_K\_M）\audit{14}，从不加载第二个模型。Ollama 0.40.7 提供模型管理和随附的 llama.cpp 库（构建号 b11738）；因为其 HTTP 接口不暴露活动序列句柄，Echo 的适配器直接加载该库。词典为 ECDICT~\cite{ecdict}。

\paragraph{增量预测（D5）。}运行时把最近一次预测的 KV 状态和排名保存在一个空闲槽中。提示变化时，复用最长公共 \emph{token} 前缀，只计算变化的尾部。完全相同的提示直接复用已存排名，不做前向。前端调度采用 25\,ms 防抖、单个在途请求、只保留最新的待处理输入。伙伴开始回复时清空这份缓存。

\paragraph{有限选项选择（D2）。}学习者以中文或英文文字看到的三种模型选择（伙伴释义、网页释义、中译英替换）都是从有限列表中选择。网页释义和中译英替换严格采用下面的形式；伙伴释义采用下文所述的 token 级变体\audit{18}。设 $P$ 为已计算的前缀，$S_i$ 为第 $i$ 项的简短任务后缀，$D_i$ 为其词典条目。每个条目获得一个标签，另加弃选标签 $0$。运行时选取
\begin{equation}
 \hat{\ell}_i = \arg\max_{\ell\in\mathcal L_i} z_\theta(\ell\mid P,S_i),\qquad
 g_i = \begin{cases}D_i[\hat{\ell}_i],&\hat{\ell}_i\ne0,\\ d_i^{\mathrm{default}}\text{（加标记）或不作改动},&\hat{\ell}_i=0,\end{cases}
\end{equation}
其中 $z_\theta$ 是限制在允许标签上的下一 token logit。$g_i$ 由词典而非生成文本提供。网页释义和中译英替换使用单 token 字母标签（A--Z 与 a--z，或 A--T），因此无论释义或短语多长，一次输出决定就够了。伙伴释义则通过 token 前缀树把贪心解码约束在中文义项字符串上，在第一个能唯一确定义项的 token 处停止（至多 16 次决定、640 token 的任务）；它每一步取最优的下一 token，选出的未必是整体得分最高的义项。开发过程中，字母标签在这条路径上选错过 \emph{dishes} 的义项。logit 不是校准过的概率。提示预填仍然耗时，因此「一次决定」绝不意味着每个请求只做一次 GPU 调用。

\paragraph{共享上下文。}伙伴释义任务必须恰好看到到目标词为止的对话，而伙伴还在继续生成。因此运行时分叉出一个短命的分支，引用伙伴\emph{实际已计算}的、直到该词的 KV 前缀（\code{llama\_memory\_seq\_cp}）。若该词结束于某个 token 内部，分支从前一个有效 token 边界开始，只重算剩余部分。随后分支用模型的模板 token 结束该回合，追加任务，作出决定，并立即释放（图~\ref{fig:architecture}）。任一时刻至多存在一个这样的分支，待处理工作限定为八条轻量描述；溢出的退回默认。分支预填每 64 个 token 让出一次，使伙伴持续流式输出。原生加锁保护伙伴待用的采样结果和 logits。网页文字作为输入只预填一次，两个针对具体出现位置的分支在同一个原生批次中共享该前缀。分支后缀每个限 640 token，合计限 704。每个分支和源序列在成功、失败或取消时都会释放。释放使池空间可以复用，但不一定减少预分配的显存。

\begin{figure}[ht]
\centering
\begin{tikzpicture}[font=\small,>=Stealth,
 box/.style={draw,rounded corners=2pt,align=center,minimum height=9mm,text width=47mm},
 base/.style={box,fill=black!5,text width=62mm}]
\node[base] (a) at (0,0) {伙伴回复 A\\历史 + 已生成到目标词的文字};
\node[box] (ac) at (7.2,0.55) {A：继续英文回复};
\node[box] (b) at (7.2,-0.55) {B：任务 + 词典义项\\决定、发布、释放};
\draw[->] (a.east) -- (ac.west);
\draw[->] (a.east) -- (b.west);
\node[base] (p) at (0,-2.0) {网页源 P\\先前段落 + 当前段落};
\node[box] (b1) at (7.2,-1.65) {B1：针对出现位置的后缀};
\node[box] (b2) at (7.2,-2.65) {B2：针对出现位置的后缀};
\draw[->] (p.east) -- (b1.west);
\draw[->] (p.east) -- (b2.west);
\node[align=center,font=\scriptsize] at (0,-3.0) {源序列依次接纳。一个模型；一个有界 KV 池。};
\end{tikzpicture}
\caption{释义分支共享其源序列已计算的 KV 前缀。伙伴释义分支读不到回复后面的文字；网页分支可以读到整个当前段落以及同一网页的先前段落。中译英替换使用一个短命的源序列和分支。}
\label{fig:architecture}
\end{figure}

\section{评测：下一输入预测是否快到可用？}
\label{sec:results}
这一范式成立与否取决于预测速度。候选列表只有反映学习者刚输入的内容时才有用。如果它在下一次击键之后才到达，描述的就是已不存在的文字，而阅读它仍要花注意力~\cite{quinn}。因此，我们在相同的模型、输入和硬件上，把预测学习者下一输入的几种方式与 HCI 文献中的延迟上限作比较。

\subsection{延迟预算}
经典指导认为，约 0.1\,s 是让响应感觉即时的上限，约 1\,s 是保持用户思路连贯的上限~\cite{nielsen,miller,card}。对打字而言，相关的时钟是击键间隔。168,000 名志愿者抄打给定句子，共 1.36 亿次击键，平均间隔为 239\,ms；最快的 10\% 参与者约 122\,ms，最慢的 10\% 约 481\,ms~\cite{dhakal}\audit{11}。组织消息的第二语言学习者可能更接近慢的一端，停顿也更多，但在一个词内部每秒仍能打出好几个字符。因此作为启发式标准，我们把在 239\,ms 内就绪的预测视为对该次击键可用，把在 100\,ms 内就绪视为舒适可用。超过 481\,ms，即使对慢速打字者也已过时。

\subsection{比较的方法}
三种方法都使用相同的 \code{qwen3:8b} 权重、相同的伙伴系统指令和相同的合成历史：同一组学习者--伙伴问答重复 12 次（24 条消息）\audit{6}。每种方法都像 Echo 一样每次输入变化发一个请求，不并行扇出。
\begin{itemize}\setlength{\itemsep}{1pt}
\item \textbf{Echo（一次前向，共享 KV）。}学习者的回合保持开放，从一次前向读出下一 token 排名，并复用上一次击键的 KV 状态（第~\ref{sec:implementation}~节）。排名随后经词典过滤。
\item \textbf{经 API 的一次前向。}相同的表述，经 Echo 随附的 Ollama 引擎的 HTTP API：生成一个 token 并返回其前 20 个对数概率。启用了 Ollama 自带的提示缓存。
\item \textbf{经 API 的提示词工程。}这是 LLM 应用索取建议的一种直截了当的做法。在同一对话中发一个聊天请求，要求模型给出\emph{「学习者最可能接着输入的 5 个词」}（或半词的补全），并用 JSON schema 把回复约束为五个字符串。
\end{itemize}
主要工作负载是在这段历史之后逐字符输入一条学习者消息 \emph{``I would like to visit a quiet park with my friends.''}：共 51 次输入变化，重复五遍。重复的伙伴句（\emph{``I like visiting quiet parks with friends\ldots''}）与键入的消息内容重合，因此这段历史比自然对话更容易续写；这对延迟影响不大，但有利于候选质量\audit{6}。另外还运行了早先记录中的四个较短的逐词序列。一个候选若是单个词典词（或四个标点之一）且与已输入的半词相符，即计为有效。这只检查形式，不检查学习者接下来是否真的会打这个词。因此我们另外记录 hit@$k$：学习者正在打或接下来要打的那个词，是否出现在前 $k$ 个候选中\audit{8}。两种一次前向方法的排名本来就经过 Echo 的词典过滤和词频补齐，其候选在很大程度上是按构造即有效；提示词工程的回答则按原样计分\audit{8}。计时在回环上从客户端测得。基准直接调用 Echo 的预测函数，因此计时不含 Echo 自身的 \code{/api/suggest} HTTP 层、25\,ms 防抖和渲染\audit{20}。两种条件于 2026 年 10 月 28 日在同一台笔记本上先后运行（NVIDIA GeForce RTX 4080 Laptop GPU，12\,GB 显存，驱动 610.62，32\,GB 内存，Windows 10），因为 12\,GB 只容得下一份模型。期间没有加载其他模型，桌面程序约占 1.3\,GB 显存\audit{14}。

\begin{table}[ht]
\centering\footnotesize\setlength{\tabcolsep}{4pt}
\begin{tabular}{@{}lrrrrrr@{}}
\toprule
 & \multicolumn{3}{c}{逐击键（250 个请求）} & 逐词 & 有效数 & 命中\\
\cmidrule(lr){2-4}
方法 & p50 / p95 & $\le$100\,ms & $\le$239\,ms & p50 范围 & （共 5） & @5\\
\midrule
Echo：一次前向，共享 KV & \textbf{5.3} / 22.4\,ms & 100\% & 100\% & 4--20\,ms & 4.76 & 94\%\\
经 API 的一次前向 & 26.3 / 34.2\,ms & 100\% & 100\% & 24--27\,ms & 4.76 & 94\%\\
经 API 的提示词工程 & 358 / 640\,ms & 0\% & 0\% & 331--444\,ms & 1.71 & 48\%\\
\bottomrule
\end{tabular}
\caption{下一输入预测，相同模型、输入和机器（2026 年 10 月 28 日）。击键各列不含每遍重复的第一个请求。「有效数」是逐击键负载上五个候选中可用候选的平均数；提示词工程方法在逐词序列上为 2.5--3.2。「命中@5」是 250 次击键中实际打出的词位于五个候选之内的比例（命中@1：88\%、88\%、40\%）\audit{8}。提示词工程的请求只有 84\%\audit{16} 在慢速打字者的 481\,ms 间隔内完成。}
\label{tab:prediction}
\end{table}

\paragraph{可用与不可用。}提示词工程方法从未在平均击键的 239\,ms 内返回，中位耗时 358\,ms。等它的列表出现时，学习者已经又打了一两个字符。它给出的词也常与已输入内容不符：5 个中平均只有 1.7 个有效，只有 48\% 的击键包含学习者实际打出的词。记录下来的回答说明了原因：被要求「预测接下来的词」时，聊天模型常常回答整个短语（\emph{``I would like to visit Paris''}）、重复已输入的内容（在 \emph{``I''} 之后给出 \emph{``I like''}），或给出词的碎片（在 \emph{``frie''} 之后给出 \emph{``nds''}）\audit{9}。并行请求候选也无济于事，因为每个请求仍要生成一段回复。两种一次前向方法在 100\,ms 内返回了全部 250 个击键请求，平均有效候选 4.76 个\audit{2}，94\% 的击键中实际打出的词在五个候选之内（88\% 是第一个候选）。这个命中率偏乐观：每条消息的 50 次输入中有 40 次落在词的内部，前缀过滤后剩下的选择很少，而且历史与键入的消息内容重合。在消息的十个词首处（第一遍重复），实际的词有七次出现在 Echo 的五个候选中，有五次出现在提示词工程的回答中\audit{8}。两种一次前向方法读取的是同一份排名：在各序列第一遍记录的全部 70 次输入中，它们的第一个候选每次都相同，56 次列表完全相同，其余只在第 3--5 位的近似并列候选上顺序或选择不同\audit{9}。因此，决定性的一步是预测本身的表述方式：在学习者自己开放的回合上读取模型的下一 token 分布，用一次前向取代了生成一段回答。

\paragraph{Echo 的余量从何而来。}Echo 中位 5\,ms、API 中位 26\,ms，这一差距有两个来源，逐击键数据可以把它们分开\audit{4}。每遍重复测得的 50 次击键中，有 39 次（78\%）发生在一个词内部。由于半截词从不送入模型，提示没有变化，Echo 直接返回已存排名，不调用模型（中位 5.0\,ms）。API 客户端同样可以缓存未变提示的结果；这里的 API 条件没有这样做，耗时 26.5\,ms。其余 11 次击键（22\%）开始一个新词，需要做一次前向：Echo 复用上一次击键的 KV 状态，中位 20.5\,ms；经 Ollama 为 26.0\,ms，而 Ollama 自带的提示缓存同样是开启的。因此原生运行时的 KV 复用省下了每次前向约五分之一的时间；总差距的大部分来自词内不重新排名。两种方法在这块 GPU 上都已低于感知阈值，所以额外的余量以其他方式起作用：它让 GPU 腾出来处理伙伴的流式回复和释义分支；它为较慢的硬件和较长的对话留出余量；而且一旦付出了上下文变化后的第一次请求，每次击键的成本就几乎与历史长度无关。在 Echo 中，24 条消息的历史下这第一次请求耗时 213--220\,ms。实际使用时，它发生在伙伴回复结束时，因为开始回复会清空预测缓存\audit{2}。API 的第一次请求不作比较，因为各遍重复之间没有重置 Ollama 的提示缓存。

\subsection{释义与中译英查询}
对阅读侧和反向查询做了同样的比较（表~\ref{tab:comparison}）。API 基线流式输出回复，每个多义词一完成就放入队列；由单个 worker 经 Ollama 的第二个槽逐个发送结构化请求。队列中已有八个词在等待时，新到的词会被丢弃且不计数；运行时则把这种溢出计为带标记的退回。各次运行都没有报告退回\audit{12}。运行时则从回复已计算的前缀分叉出分支。差别主要体现在释义\emph{何时}出现。经 API 时，最后一个释义在英文全文完成后约 4\,s 才到，学习者会在没有帮助的情况下读完大半条回复。使用运行时，全部释义在请求后 2.4\,s 上屏，其中几个早于正文结束。经 API 的自由翻译很快但没有边界：14 个样例中只有 5 个的回答等于某个词典义项\audit{7}，有些是解释，或翻译了错误的词（\emph{dishes}$\to$洗碗，即 ``wash dishes''），而且没有一个弃选。运行时与结构化 API 路径之间的准确率差距主要来自提示和弃选策略，而非 KV 复用。这些样例也在提示开发中使用过。中译英查询是单次、由用户触发的请求，这里 API 路径更快，因为运行时必须新建源序列和分支。前缀共享在此没有好处。

\begin{table}[ht]
\centering\footnotesize\setlength{\tabcolsep}{4pt}
\begin{tabular}{@{}llrrr@{}}
\toprule
任务 & 指标 & 运行时 & API（受限） & API（自由）\\
\midrule
多义词释义 & 延迟 p50 / p95 & 204 / 267\,ms & 717 / 817\,ms & 86 / 263\,ms\\
 & 答案可接受\audit{7} & 14/14 & 10/14 & 8/14\\
 & 输出为词典义项 & 14/14 & 14/14 & 5/14\\
\addlinespace
伙伴回复加释义 & 英文全文完成 & 1.75\,s & 1.34\,s & ---\\
 & 全部释义上屏 & 2.37\,s & 5.38\,s & ---\\
 & 正文结束前到达的释义 & 3--5 / 6--10 & 0 / 5--7 & ---\\
\addlinespace
中译英 & 延迟 p50 / 可接受 & 107\,ms / 5/5 & 43\,ms / 5/5 & 65\,ms / 5/5\\
\bottomrule
\end{tabular}
\caption{释义与反向查询，相同模型和机器（2026 年 10 月 28 日）。准确率按 14 个释义样例和 5 个反查词计。每项运行三次，判定结果完全相同（自由回答有两例措辞不同），因此重复只用于延迟（分别计时 42 和 15 次）；另计时三条回复。三个释义样例期望弃选：两个左侧语境不足，一个所需义项不在词典中\audit{5}。剔除它们后，计数为 11/11、10/11 和 8/11。自由回答按宽松标准判定，且无法弃选。}
\label{tab:comparison}
\end{table}

\paragraph{注意事项。}每种条件只运行一次，在一台机器上，有轻度后台 GPU 负载，只用一个模型。回复长度随采样而变化。云端 API 会给每次击键加上网络往返，还会把学习者的文字发到设备之外，Echo 在设计上避免了这一点。这些阈值只界定感知延迟。及时的候选是否帮助学习者，是第~\ref{sec:limitations}~节的问题。

\subsection{预测缓存与网页释义}
其余测量来自仓库的验收脚本，在同一天、同一台机器、同一模型上运行\audit{15}。预测方面，生命周期脚本在六个追加或改动的输入上比较缓存排名与冷计算排名：每个缓存请求复用了提示中的 143--144 个 token，只计算一到三个。缓存与冷计算的 GPU 分数略有差异（共同前 20 项的对数概率最大差 0.21），而相同提示的排名复用是精确的。网页释义方面，四段手写短文各含两个目标出现位置：河岸与金融的 \emph{bank}、重量与照明的 \emph{light}、声音与结构的 \emph{sound}、预订与印刷品的 \emph{book}。测试先运行单分支模式，再运行双分支模式，共对八个独立出现位置作 16 次目标决定。全部通过了可接受义项检查，每次只用一次输出决定（表~\ref{tab:reader}）。

\begin{table}[ht]
\centering\small
\begin{tabular}{@{}lrr@{}}
\toprule
短文目标 & 单分支（ms） & 至多两分支（ms）\\
\midrule
bank：岸 / 银行 & 493.8 & 184.0\\
light：轻的 / 光源 & 696.7 & 557.2\\
sound：声音 / 健全的 & 423.0 & 381.5\\
book：预订 / 书 & 249.2 & 220.5\\
\bottomrule
\end{tabular}
\caption{一次网页运行（2026 年 10 月 28 日）：完成每段短文中全部合格释义所需时间，不只是两个目标。初始词典事件为 1.0--6.0\,ms；输入前缀为 131--133 个 token。该测试单独为每段短文选义，不使用第~\ref{sec:design}~节所述的网页上下文。测试按顺序执行，存在冷/热状态混杂，因此这些耗时不能证明批处理带来加速。}
\label{tab:reader}
\end{table}

\subsection{生命周期与隔离}
1,000 次分叉/释放检查用时 26.12\,s\audit{15}，结束时活动的源序列和分支句柄均为零；分叉与释放计数相等（包括同一服务会话中的其他检查，各 1,132 次）。在固定种子（73）下，插入或不插入一个短的或长的释义分支，伙伴回复文字完全一致。同一次运行中，回复的 token 间隔 P50 单独时为 14.5--17.0\,ms，插入长分支后 P95 从 17.7\,ms 升至 40.0\,ms，所以释义对伙伴的流式输出并非没有代价。这些检查支持生命周期计数和有限的采样隔离，并不测量物理 KV 字节，也不证明在所有负载下都无泄漏或确定。

\section{相关工作}
\label{sec:related}
\paragraph{语言模型引导的文本输入。}Dasher 让人通过在按语言模型预测比例分配的空间中驾驭方向来书写~\cite{dasher}，是把人当作解码器这一思路最接近的概念先例。SpeakFaster 利用大语言模型和对话上下文，为患 ALS 的眼动打字者扩展缩写输入~\cite{speakfaster}。两者都着眼于已经知道自己要说什么的人的输入速度或无障碍。Echo 把引导的思路用于仍在用第二语言组织话语的学习者：其预测来自对话伙伴自己的上下文，并在每一步提供词汇支持。

\paragraph{响应时间。}自 Miller 对对话式事务的分析以来，交互系统的响应时间上限一直很稳定~\cite{miller}：约 0.1\,s 感觉即时，约 1\,s 保持思路连贯~\cite{card,nielsen}。大规模打字数据把平均击键间隔定在 239\,ms 附近~\cite{dhakal}。Quinn 和 Zhai 表明，建议能省击键，但关注建议的成本可能使平均输入变慢~\cite{quinn}。我们据此推断，迟到或错误的建议只增加这份成本而没有收益\audit{10}。Echo 把这些数字作为设计预算（第~\ref{sec:results}~节）。

\paragraph{建议对写作的影响。}预测性建议会改变人们写下的内容。Arnold 等发现，借助建议写的图片说明更短、更可预测~\cite{predictable}。Buschek 等发现，多条并行的短语建议有助于构思但降低效率，且非母语写作者似乎从更多建议中获益更多~\cite{parallel}。CoAuthor 大规模记录了人机协同写作会话~\cite{coauthor}。这些结果既促成了 Echo 的设计（提供若干单词选项，并允许学习者随时自由输入），也构成其核心的未决风险：把学习者引向可预测的语言。

\paragraph{支撑技术。}基于释义的词义消歧已很成熟~\cite{glossbert}，词汇受限解码和语法引导解码~\cite{hokamp,outlines}、服务系统中的共享前缀 KV 复用~\cite{vllm,sglang,hydragen}也是如此。Echo 使用这些技术，并不主张推进它们。它的特定需求是本地的：在一块消费级 GPU 上，释义分支必须从仍在生成的回复中、在过去的位置分叉，且不明显拖慢对话。Echo 的分支在一次选择后即丢弃，并不实现思维树（Tree of Thoughts）~\cite{tot} 那样的推理路径搜索。

\section{局限、假设与评估计划}
\label{sec:limitations}
\paragraph{尚无学习者证据。}本报告记录了一个可运行的系统，并未表明学习者写得更好、学到了词汇、更有信心，或相比普通预测键盘、词典、整条回复建议更偏好 Echo。「已会」标记记录的是学习者的选择，而不是测得的掌握程度。

\paragraph{假设。}该设计引出若干可检验的问题：
\textbf{H1}，与自由输入相比，学习者引导的单词候选让学习者完成更多、更长的回复，同时比整条回复建议保留更多源自学习者的内容；
\textbf{H2}，依照可预测性效应~\cite{predictable}，使用候选会降低词汇多样性，且这一效应取决于显示的候选数量~\cite{parallel}；
\textbf{H3}，带释义的候选和就地中译英替换让学习者接触新词，之后对这些词的识别率高于基线；
\textbf{H4}，可见的退回标记帮助学习者判断释义何时不确定。

\paragraph{评估计划。}可以让两个水平层次的中文母语学习者参加被试内研究，在相同的对话任务中比较三种输入方式：带词典的自由输入、Echo、整条回复建议。产出指标为完成度、长度以及自行输入词与接受词的比例；语言指标包括词汇多样性、错误率以及与模型 top-1 续写的相似度；体验指标包括感知自主和工作负荷。对首次通过候选或替换遇到的词做延迟识别测试，可检验 H3。记录的交互轨迹可支持建议层面的分析。

\paragraph{系统局限。}释义质量受词典限制（词典可能缺少短语义项或义项划分不佳），也受模型未经校准的选择限制。伙伴释义刻意保持因果性：在 ``the bank'' 处，后文本可以确定含义，但分支看不到。网页和替换的标签把裸 \code{0} 与带前导空格的字母 token 配对；项目笔记指出这种分词不对称可能造成标签偏差，尚未测量。16 次网页决定只是冒烟测试，不是准确率估计。候选过滤无法验证词边界，对话模板按模型族编写（Qwen 和 Gemma），KV 池和队列小且固定，中文输入桥仅适用于 Windows。

\section{代码与数据}
\label{sec:reproduction}
源代码公开于 \url{https://github.com/cyx3705/2026-040}。本报告描述的系统是提交 \code{143bec4} 的应用代码\audit{13}。基准脚本以及第~\ref{sec:results}~节每个数字背后的原始 JSON 输出，位于仓库的 \code{technical-report} 分支，分别在 \code{b-Code-Echo/tests/} 和 \code{b-Office/technical-report/data/} 下\audit{15}。全部测量于 2026 年 10 月 28 日完成。表~\ref{tab:prediction} 和表~\ref{tab:comparison} 来自 \code{api\_comparison.py}，两种模式各运行一次；表~\ref{tab:reader} 来自 \code{native\_reader\_acceptance.py}；第~\ref{sec:results}~节的缓存对比和生命周期检查来自 \code{native\_acceptance.py}。图~\ref{fig:walk} 的候选列表最初记录于 2026 年 10 月 23 日，在 10 月 28 日的输出中原样复现。本报告伴随实现提供，并不声称是公开发布版或不可变的基准包；重新运行脚本得到的是新的观察，而不是完全相同的轨迹。

原生环境需要 Windows x64、Python 3.12 或更高版本、Visual Studio C++ Build Tools 和 Windows SDK、Ollama 0.40.7 及本地 \code{qwen3:8b} 模型，以及 ECDICT SQLite 词典\audit{14}。在依赖就绪、Echo 空闲运行时，以下在仓库根目录执行的命令会调用相关检查。GPU 检查应依次运行。
\begin{quote}\footnotesize\ttfamily
uv run -{}-project b-Code-Echo python\\
\hspace*{1em}b-Code-Echo/tests/native\_reader\_acceptance.py\\[3pt]
uv run -{}-project b-Code-Echo python\\
\hspace*{1em}b-Code-Echo/tests/native\_acceptance.py -{}-cycles 1000\\[3pt]
uv run -{}-project b-Code-Echo python\\
\hspace*{1em}b-Code-Echo/tests/api\_comparison.py -{}-mode native -{}-out native.json
\end{quote}
上面每两行为一条命令。对比脚本的 API 模式需在 Echo 关闭、且 \code{ECHO\_KV\_ADAPTER\_URL} 设为空字符串时运行；它会自行启动和停止私有 Ollama 引擎。网页测试在内存中临时把目标词元视为未学会，不写入学习状态或设置。其他已学会的词会被跳过，因此本地学习状态的差异会改变总工作量和耗时。

默认情况下，一切都在学习者的机器上运行。扩展把提取出的段落文字连同网页的源和路径（或视频 ID）作为上下文键发给 Echo 服务，从不发送页面 HTML、Cookie 或表单内容。Echo 把第~\ref{sec:design}~节所述的网页上下文存入数据目录，以该键的摘要命名。Echo 也可以选择通过访问密钥共享给私有网络中的其他设备，此时段落文字会发往那台主机\audit{13}。扩展排除可编辑区域、表单和代码区域。链接需 Alt+点击才切换已会状态，因此正常导航不受影响。这些是实现边界，不是安全审计。把网页文字作为数据加引号可以减少指令混淆，但不能证明对对抗性内容免疫。源码许可以仓库中的声明为准。

\begingroup
\footnotesize
\begin{thebibliography}{99}
\setlength{\itemsep}{2pt}
\bibitem{miller} R. B. Miller. Response time in man-computer conversational transactions. In \emph{Proc.\ AFIPS Fall Joint Computer Conference}, 1968. doi:10.1145/1476589.1476628.
\bibitem{card} S. K. Card, G. G. Robertson, and J. D. Mackinlay. The information visualizer, an information workspace. In \emph{Proc.\ CHI}, 1991. doi:10.1145/108844.108874.
\bibitem{nielsen} J. Nielsen. \emph{Usability Engineering}, chapter 5. Academic Press, 1993.
\bibitem{dhakal} V. Dhakal, A. M. Feit, P. O. Kristensson, and A. Oulasvirta. Observations on typing from 136 million keystrokes. In \emph{Proc.\ CHI}, 2018. doi:10.1145/3173574.3174220.
\bibitem{quinn} P. Quinn and S. Zhai. A cost-benefit study of text entry suggestion interaction. In \emph{Proc.\ CHI}, 2016. doi:10.1145/2858036.2858305.
\bibitem{dasher} D. J. Ward, A. F. Blackwell, and D. J. C. MacKay. Dasher---a data entry interface using continuous gestures and language models. In \emph{Proc.\ UIST}, 2000. doi:10.1145/354401.354427.
\bibitem{speakfaster} S. Cai et al. Using large language models to accelerate communication for eye gaze typing users with ALS. \emph{Nature Communications}, 15:9449, 2024. doi:10.1038/s41467-024-53873-3. Preprint arXiv:2312.01532.
\bibitem{predictable} K. C. Arnold, K. Chauncey, and K. Z. Gajos. Predictive text encourages predictable writing. In \emph{Proc.\ IUI}, 2020. doi:10.1145/3377325.3377523.
\bibitem{parallel} D. Buschek, M. Z\"urn, and M. Eiband. The impact of multiple parallel phrase suggestions on email input and composition behaviour of native and non-native English writers. In \emph{Proc.\ CHI}, 2021. doi:10.1145/3411764.3445372.
\bibitem{coauthor} M. Lee, P. Liang, and Q. Yang. CoAuthor: Designing a human-AI collaborative writing dataset for exploring language model capabilities. In \emph{Proc.\ CHI}, 2022. doi:10.1145/3491102.3502030. Preprint arXiv:2201.06796.\audit{21}
\bibitem{glossbert} L. Huang, C. Sun, X. Qiu, and X. Huang. GlossBERT: BERT for word sense disambiguation with gloss knowledge. In \emph{Proc.\ EMNLP-IJCNLP}, 2019. arXiv:1908.07245.
\bibitem{hokamp} C. Hokamp and Q. Liu. Lexically constrained decoding for sequence generation using Grid Beam Search. In \emph{Proc.\ ACL}, 2017. doi:10.18653/v1/P17-1141.
\bibitem{outlines} B. T. Willard and R. Louf. Efficient guided generation for large language models. 2023. arXiv:2307.09702.
\bibitem{vllm} W. Kwon et al. Efficient memory management for large language model serving with PagedAttention. In \emph{Proc.\ SOSP}, 2023. arXiv:2309.06180.
\bibitem{sglang} L. Zheng et al. SGLang: Efficient execution of structured language model programs. In \emph{Proc.\ NeurIPS}, 2024. arXiv:2312.07104.
\bibitem{hydragen} J. Juravsky et al. Hydragen: High-throughput LLM inference with shared prefixes. 2024. arXiv:2402.05099.
\bibitem{tot} S. Yao et al. Tree of Thoughts: Deliberate problem solving with large language models. In \emph{Proc.\ NeurIPS}, 2023. arXiv:2305.10601.
\bibitem{ecdict} skywind3000. ECDICT. Software and English--Chinese dictionary. \url{https://github.com/skywind3000/ECDICT}. Accessed October 2026.
\bibitem{llamacpp} llama.cpp contributors. llama.cpp. Software; build b11738 as bundled with Ollama 0.40.7. \url{https://github.com/ggml-org/llama.cpp}. Accessed October 2026.
\end{thebibliography}
\endgroup

\clearpage
\section*{附：核验意见处理情况（2026-10-28）}
\small
问题依据与数据拆分见 \code{REPORT-001-核验-2026-10-28.md}。「已改」表示问题已在正文中消除。另按作者要求删除了测试通过情况、浏览器验证范围、终端环境等运维记录。
\begin{enumerate}\setlength{\itemsep}{1pt}
\item[A1] 已改：用修订稿重新运行 \code{reportBuild}，英文 PDF（11 页）与源码包均已更新，扩展测试路径已是新位置。
\item[A2] 已改：摘要和第 5.2 节改为「除首次请求外全部 $<$100\,ms」，写明首次请求约 214\,ms 及其发生时机，并写明平均有效 4.76 个。
\item[A3] 已改（仅披露）：第 3.1 节写明 Shift+W/S 在候选有效时会拦截大写 W/S。代码未改，键位冲突仍然存在。
\item[A4] 已改：第 5.2 节改写为「Echo 的余量从何而来」，分开报告词内复用（78\%，4.8\,ms）和需要前向（22\%，20.1 对 27.6\,ms）；摘要与贡献列表同步修改。
\item[A5] 已改：表 2 题注写明两个样例左侧语境不足，一个是词典缺少义项。
\item[A6] 已改：写明历史为同一组问答重复 12 次（24 条消息），且与键入句内容重合。
\item[A7] 已改：表 2 改按独立样例计（14/14、10/14、8/14；14/14、14/14、4/14；5/5），重复次数只用于延迟。
\item[B8] 已改：写明「有效」只检查形式及两组的不对等；10-28 重跑时保存了候选词，并新增命中@1/@5（Echo 88\%/94\%，提示词工程 40\%/48\%）和词首处的拆分。
\item[B9] 已改：两处说法改为依据记录数据——提示词工程的实际回答举例；两种一次前向方法第一个候选 70/70 相同、完整列表 56/70 相同。
\item[B10] 已改：相关工作中写明「比没有更糟」是本文的推论。
\item[B11] 已改：写明最快/最慢各 10\%、抄打任务，239\,ms 是启发式阈值。
\item[B12] 已改：写明单 worker 串行发送，以及队列满时丢弃且不计数的不对称。
\item[B13] 已改：第 4、8 节固定为提交 \code{143bec4} 并给出公开仓库地址；隐私说明按该提交更新（上下文键含网页源和路径、网页上下文落盘、可选共享）。
\item[B14] 已改：补充 RTX 4080 Laptop 12\,GB、驱动 610.62、32\,GB 内存、Windows 10、Ollama 0.40.7、llama.cpp b11738、Q4\_K\_M。
\item[B15] 已改：10-28 重跑全部测量（含 1000 次生命周期 26.12\,s 和表 3），原始输出归档在 \code{data/2026-10-28/}；无法复核的早先记录已删除。
\item[C16] 已改：「只有 84\%」。
\item[C17] 已改：图 1 改为记录顺序 maker, makeup / making, makeshift。
\item[C18] 已改：式 (1) 后补充 token 与词的说明；式 (2) 注明只严格适用于单 token 标签，并注明 trie 解码是逐步贪心。
\item[C19] 已改：依赖 \code{/api/suggest} 端到端测量的早先记录已随重测删除，口径问题不再存在。
\item[C20] 已改：第 5.2 节注明计时不含 \code{/api/suggest} 这一层 HTTP。
\item[C21] 已改：CoAuthor 补 DOI；ECDICT 与 llama.cpp 补访问日期，llama.cpp 补构建号。GlossBERT 的 ACL 编号未核实，未补。
\end{enumerate}
\end{document}
